\chapter{Статусная заморозка дискретной компакттонной ветви}
\label{ch:v6-spectral-transition-discrete-compacton-branch-status-freeze-gate}

\section{Задача заморозки}

Компакттонная ветвь прошла полный цикл проверок: существование,
флоке-устойчивость, физический масштаб, динамический захват, аффинный
$C_4$-селектор, граничную эта-фазу, энергетическое вырождение и
характерно-разрешённую радиацию. Настоящая глава не вводит новый член и
не повторяет поиск диссипации. Она сопоставляет восемь машинных
сертификатов с контрактом R0--R6 и фиксирует границу применимости
полученной модели.

Такое разделение существенно. Точные compact-support решения нелинейных
решёток и локализованные состояния квантовых блужданий являются
самостоятельными математическими объектами, но их существование не
доказывает динамическое образование физической частицы
\cite{compacton-freeze-lee2016,compacton-freeze-bisio2021}. Аналогично,
ненулевая связь с непрерывным каналом сама по себе не задаёт
стабилизирующее затухание: необходимо проверить знак потока и истощение
дискретной моды \cite{compacton-freeze-fano1961}.

\section{Сохранённое точное ядро}

После заморозки сохраняются семь результатов.
\begin{enumerate}
 \item На составном хиральном носителе существуют точные двухузловые
 compacton-решения
 \begin{equation}
  \kappa=2(2m+1)\pi,
  \qquad \kappa_{\min}=2\pi,
  \qquad F(\Psi_\pm)=\pm i\Psi_\pm.
 \end{equation}
 \item Одношаговые ветви $\pm i$ лежат внутри непрерывного точного
 двухузлового многообразия с законом
 \begin{equation}
  F^2(\Psi)=-\Psi.
 \end{equation}
 \item На проверенных периодических объёмах редуцированная
 четырёхшаговая монодромия не имеет расширяющих мультипликаторов.
 \item На характерном пространстве существует безразмерный дефект
 \begin{equation}
  D_\chi=4w_+w_-,
 \end{equation}
 нули которого равны паре $\pm i$.
 \item После явного выбора оси граничное условие может согласованно
 нарушить $S_4$ до одной подгруппы $C_4$.
 \item Для эта-ориентированной квадратуры вычислены точные величины
 \begin{equation}
  \rho_{\rm rad}(0)=2\pi,
  \qquad \Gamma_{\rm cycle}=4\pi^2|\delta|^2.
 \end{equation}
 \item Полученная система является воспроизводимым стендом для проверки
 нелинейной локализации, ненормальных переходных процессов и
 конечнопорогового радиационного ухода.
\end{enumerate}

Это положительное ядро не отменяется отрицательным физическим вердиктом.
Оно только получает более узкую интерпретацию.

\section{Сводка по контракту R0--R6}

\begin{center}
\begin{tabular}{p{0.08\textwidth}p{0.19\textwidth}p{0.58\textwidth}}
\toprule
Тест & Статус & Причина \\
\midrule
R0 & частично & Хиральный носитель задан, но ось и граница
$S_4\to C_4$ выбраны условно, а не выведены из полного родителя
$\mathcal H_{15}$. \\
R1 & частично & Нелинейное унитарное блуждание едино для compacton,
но граница, аффинный селектор и открытая динамика не следуют из одного
Real/gauge-функционала. \\
R2 & провал & Точное многообразие не является изолированным аттрактором;
в проспективном протоколе получено $0/36$ захватов. \\
R3 & провал & Эта/Pfaffian-фаза даёт ориентацию, но нулевую скорость.
Коэффициент $4\pi^2$ описывает излучение ядра, а не вероятность рождения
или захвата. \\
R4 & провал & Имеется локально нейтральное решение, но сохраняются
непрерывное многообразие, нейтральная квадратура и конечнопороговый уход;
единственного устойчивого endpoint нет. \\
R5 & провал & Из $\kappa=2\pi$ следует лишь
$EL=\pi\hbar c$; абсолютные масса, размер, время жизни и скорость не
предсказаны. \\
R6 & пройден & Все положительные и отрицательные утверждения имеют
воспроизводимые JSON-сертификаты и заранее заданные стоп-критерии. \\
\bottomrule
\end{tabular}
\end{center}

Итоговый счёт равен одному пройденному, двум частичным и четырём
проваленным пунктам. Согласно контракту допуска, частичный успех R0 и R1
не компенсирует провалы R2--R5.

\section{Почему найденный форм-фактор не спасает ветвь}

Последний тест устранил последнюю численную неопределённость. Аффинная
линия характера $-1$ действительно имеет ненулевой форм-фактор в
массовый безщелевой континуум сдвига. Однако линейная проекция источника
остаётся единичной, последовательные исходящие пакеты ортогональны, а
полная нелинейная эволюция уменьшает вероятность всего двухузлового ядра.
При этом нежелательный характер не затухает и одна поперечная квадратура
остаётся точной касательной compacton-многообразия.

Следовательно, число $4\pi^2$ имеет правильное происхождение, но
неправильный физический смысл. Это коэффициент радиационной
неустойчивости, а не выведенный $\gamma$ амплитудного захвата.

\begin{nogocriterion}[Запрет чтения излучения как рождения]
Нельзя объявлять ветвь открытой системой с захватом только потому, что
найдена ненулевая спектральная плотность. Необходимы одновременно
истощение нежелательной характерной амплитуды, отрицательная действительная
скорость всех физических поперечных квадратур и сохранение локализованного
ядра. В текущем родителе не выполнено ни одно из трёх совместных условий.
\end{nogocriterion}

\section{Стоп-правило и условия переоткрытия}

После настоящей главы запрещается продолжать серию добавлением ещё одной
фазы, ещё одного скалярного граничного коэффициента или феноменологического
Lindblad-оператора. Переоткрытие допустимо только если один новый
родитель одновременно выводит:
\begin{enumerate}
 \item физический масштаб решётки или локализации;
 \item ненулевой бассейн захвата из общих или вакуумно совместимых
 начальных состояний;
 \item отрицательное действительное затухание всех физических
 поперечных характерных квадратур при сохранении ядра;
 \item граничный $C_4$-смеситель или соответствующий открытый канал;
 \item нормированную физическую скорость либо слепой observable.
\end{enumerate}
Эти условия должны следовать из одной меры и одного Real/gauge-родителя.
Выполнение лишь одного пункта не снимает заморозку.

\begin{theorem}[Статус компакттонной ветви]
Дискретная составная модель содержит точное нелинейное локализованное
решение и аналитически нормированный радиационный канал. Но она не имеет
эндогенного запуска, вероятности или скорости рождения, единственного
устойчивого endpoint и абсолютного физического масштаба. Поэтому
компакттон сохраняется как математический результат и вычислительный
стенд, но закрывается как автономный механизм рождения материи.
\end{theorem}

\section{Следующий гейт}

Следующий гейт
\path{version6_spectral_transition_post_compacton_program_reprioritization_gate}
должен вернуться к полному ledger Тома VI и выбрать следующий маршрут не
по локальной привлекательности механизма, а по тому, какой кандидат
одновременно атакует провалы R1--R5.

Нулевая гипотеза: ни одна оставшаяся подсказка текущей версии не даёт
такого совместного родителя. Стоп-критерий: если все кандидаты снова
требуют независимых масштаба, trigger и диссипации, Том VI завершается
честным отрицательным динамическим статусом без новой модели.

\begin{protocol}
\begin{enumerate}
 \item сведённых машинных гейтов: \textbf{восемь};
 \item точный двухузловой compacton: \textbf{сохранён};
 \item минимальная связь: \textbf{$2\pi$};
 \item общий закон точного многообразия: \textbf{$F^2=-1$};
 \item расширяющие редуцированные флоке-моды: \textbf{не найдены};
 \item захваты общих состояний: \textbf{$0/36$};
 \item абсолютный масштаб: \textbf{не выведен};
 \item эта/Pfaffian-скорость: \textbf{ноль};
 \item радиационный коэффициент: \textbf{$4\pi^2$};
 \item радиация стабилизирует ядро: \textbf{нет};
 \item R0--R6: \textbf{1 пройден, 2 частично, 4 провалены};
 \item автономное рождение материи: \textbf{не допущено};
 \item статус ветви: \textbf{заморожена с сохранением точного ядра}.
\end{enumerate}
\end{protocol}

Машинный сертификат сохранён в
\path{s2t_v6_spectral_transition_discrete_compacton_branch_status_freeze_gate_results.json}.

\begin{thebibliography}{9}
\bibitem{compacton-freeze-lee2016}
C.-W.~Lee, P.~Kurzyński, and H.~Nha,
\emph{Quantum Walk as a Simulator of Nonlinear Dynamics: Nonlinear Dirac
Equation and Solitons}, Physical Review A 93 (2016), 042325,
\path{arXiv:1512.08358}.

\bibitem{compacton-freeze-bisio2021}
A.~Bisio, N.~Mosco, and P.~Perinotti,
\emph{Scattering and Perturbation Theory for Discrete-Time Dynamics},
Physical Review Letters 126 (2021), 250503,
\path{arXiv:1912.09768}.

\bibitem{compacton-freeze-fano1961}
U.~Fano,
\emph{Effects of Configuration Interaction on Intensities and Phase Shifts},
Physical Review 124 (1961), 1866--1878,
\path{doi:10.1103/PhysRev.124.1866}.
\end{thebibliography}